\documentclass[11pt]{article}
\usepackage[a4paper,margin=2.4cm]{geometry}
\usepackage{kotex}
\setmainhangulfont{Noto Serif CJK KR}
\setsanshangulfont{Noto Sans CJK KR}
\usepackage{microtype}
\usepackage{graphicx}
\usepackage{booktabs}
\usepackage{tikz}
\usetikzlibrary{arrows.meta,positioning,shapes.geometric}
\usepackage{pgfplots}
\pgfplotsset{compat=1.17}
\usepackage[hidelinks]{hyperref}
\hypersetup{
  pdftitle={ResearchPDF: 논문을 위한 Chrome PDF 뷰어},
  pdfauthor={ResearchPDF},
  pdfsubject={시작 가이드에서 여는 ResearchPDF 둘러보기},
  pdfkeywords={PDF viewer, papers, annotations, Google Drive sync}
}

\title{\textbf{ResearchPDF: 논문을 위한 Chrome PDF 뷰어}}
\author{ResearchPDF 둘러보기\\[2pt]\small\url{https://research-pdf.croksuter.com}}
\date{}

\renewcommand{\refname}{참고문헌}
\renewcommand{\figurename}{그림}
\renewcommand{\tablename}{표}
\renewcommand{\abstractname}{요약}
\begin{document}
\maketitle

\begin{abstract}
\noindent
지금 이 문서는 ResearchPDF 안에서 열려 있습니다. 확장 프로그램 자체를 소개하는
짧은 논문이라, 여기서 설명하는 기능을 바로 이 자리에서 써 볼 수 있습니다. 지금
보고 있는 탭, 위쪽의 논문 정보 줄, 필기 도구, 그림 캡처가 모두 그렇습니다. 각 절은
\emph{해 보기} 한 줄로 끝납니다.
\end{abstract}

\section{논문은 탭 하나에}
웹이나 컴퓨터에서 연 PDF는 웹 탭 사이에 흩어지지 않고, 자체 탭 줄을 가진
\emph{PDF 탭} 하나에 모입니다. 논문은 \emph{프로젝트}로 나눕니다. 새로 연 PDF는 모두
기본 프로젝트에서 시작하고, \emph{프로젝트로 이동}으로 다른 곳에 보냅니다. 프로젝트는
폴더로 묶고, 끌어서 순서를 정하고, 아이콘·이모지·색을 줄 수 있습니다. 그 아이콘은
Chrome 탭 아이콘으로도 보입니다. 프로젝트 탭을 닫았다가 다시 열면 같은 논문들이
그대로 돌아옵니다.

\smallskip\noindent\emph{해 보기:} \textsf{Alt+Shift+\(\rightarrow\)}로 다음 탭,
\textsf{Alt+W}로 닫기, \textsf{Alt+Shift+T}로 다시 열기. 집 모양 버튼은 \emph{홈}을
엽니다. 열어 본 모든 문서를 읽은 정도, 필터, 검색과 함께 보여 줍니다.

\section{논문을 알아보기}
DOI, arXiv 식별자, 제목으로 문서를 논문으로 알아보면, 도구 모음 아래 줄에 게재처와
연도, 인용 수(연도별 그래프 포함), 참고문헌, 바로 복사할 수 있는 BibTeX·APA가
나옵니다. 데이터는 공개 학술 데이터베이스에서 가져오며(표~\ref{tab:sources}),
보내는 것은 식별자뿐입니다.

\begin{table}[h]
\centering
\caption{논문 정보 줄의 데이터 출처.}
\label{tab:sources}
\begin{tabular}{@{}lll@{}}
\toprule
출처 & 제공하는 것 & 참고문헌 \\
\midrule
OpenAlex & 게재처, 연도별 인용, 참고문헌 & \cite{openalex} \\
Crossref & DOI, 출판사, 참고문헌 수 & \cite{crossref} \\
Semantic Scholar & 프리프린트·출판본 통합 인용 수 & \cite{s2} \\
arXiv & 프리프린트와 버전 & --- \\
\bottomrule
\end{tabular}
\end{table}

\noindent\emph{해 보기:} arXiv나 출판사에서 아무 논문이나 열어 보세요. 이 문서는
출판된 논문이 아니라서 정보 줄에 그렇게 나옵니다.

\section{따라오는 필기}
형광펜, 펜, 텍스트 메모는 문서마다 저장되고, 문서는 파일 이름이 아니라 내용으로
알아봅니다. Google 계정을 연결하면 필기와 읽던 페이지가 \emph{내} Drive의 앱 전용
숨김 폴더를 통해 모든 Chrome 프로필과 기기를 따라옵니다(그림~\ref{fig:sync}).
PDF 파일 자체는 올리지 않습니다.

\begin{figure}[h]
\centering
\begin{tikzpicture}[
  node distance=1.2cm and 1.6cm,
  dev/.style={draw, rounded corners=3pt, minimum width=2.6cm, minimum height=0.9cm, fill=blue!6},
  cloud/.style={draw, ellipse, minimum width=3.6cm, minimum height=1.2cm, fill=green!8},
  link/.style={{Latex[length=2mm]}-{Latex[length=2mm]}, thick}
]
\node[cloud] (drive) {내 Google Drive (앱 전용 폴더)};
\node[dev, below left=of drive] (lap) {노트북};
\node[dev, below=of drive] (lab) {연구실 PC};
\node[dev, below right=of drive] (pro) {다른 프로필};
\draw[link] (lap) -- node[left, font=\small] {동기화} (drive);
\draw[link] (lab) -- (drive);
\draw[link] (pro) -- (drive);
\end{tikzpicture}
\caption{필기와 읽던 위치는 내 Drive의 앱 전용 폴더로 동기화됩니다. 개발자 서버는
거치지 않습니다.}
\label{fig:sync}
\end{figure}

\noindent\emph{해 보기:} 필기 도구(펜 버튼)를 열고 이 문단의 한 문장에 형광펜을 칠한
뒤, 다른 기기에서 이 문서를 다시 열어 보세요.

\section{키 하나로 그림 캡처}
\textsf{S}를 누르고 원하는 영역을 끌면 출처와 함께 이미지로 복사됩니다. \textsf{S}를
한 번 더 누르면 페이지의 그림과 표를 찾아 줍니다. 기기 안에서 도는 작은 레이아웃
모델\cite{ppdoclayout}과 PDF의 캡션을 함께 씁니다. 하나를 누르면 복사됩니다.
그림~\ref{fig:plot}과 표~\ref{tab:sources}로 해 보세요.

\begin{figure}[h]
\centering
\begin{tikzpicture}
\begin{axis}[
  ybar, bar width=10pt, width=0.8\linewidth, height=5cm,
  xlabel={연도}, ylabel={인용 수},
  symbolic x coords={2019,2020,2021,2022,2023,2024,2025},
  xtick=data, ymin=0, enlarge x limits=0.08,
  nodes near coords, every node near coord/.append style={font=\scriptsize}
]
\addplot[fill=blue!35, draw=blue!60] coordinates {(2019,12) (2020,41) (2021,88) (2022,143) (2023,190) (2024,231) (2025,204)};
\end{axis}
\end{tikzpicture}
\caption{논문 정보 줄이 그리는 연도별 인용 그래프의 예(예시 수치).}
\label{fig:plot}
\end{figure}

\noindent\emph{해 보기:} 지금 \textsf{S}를 두 번 누르고 그림 하나를 눌러 보세요.

\begin{thebibliography}{9}
\bibitem{openalex} J.~Priem, H.~Piwowar, and R.~Orr. OpenAlex: A fully-open
index of scholarly works, authors, venues, institutions, and concepts.
\emph{arXiv:2205.01833}, 2022.
\bibitem{crossref} G.~Hendricks, D.~Tkaczyk, J.~Lin, and P.~Feeney. Crossref:
The sustainable source of community-owned scholarly metadata.
\emph{Quantitative Science Studies}, 1(1):414--427, 2020.
\bibitem{s2} R.~Kinney et al. The Semantic Scholar Open Data Platform.
\emph{arXiv:2301.10140}, 2023.
\bibitem{ppdoclayout} T.~Sun et al. PP-DocLayout: A unified document layout
detection model to accelerate large-scale data construction.
\emph{arXiv:2503.17213}, 2025.
\end{thebibliography}

\end{document}
